\chapter{Einschränkungen, Intervalle und Konvexität}
\label{b12guide:ch:subsets}

\section{Einschränkungen}

\begin{definition}[Induzierte Halbordnung]
\label{b12guide:def:restriction}
Für $B\subseteq P$ heißt die auf $B$ eingeschränkte Relation
\[
 x\leq_B y\quad\Longleftrightarrow\quad x\leq y
 \qquad(x,y\in B)
\]
\emph{induzierte Ordnung} oder \emph{Restriktion}.
Ihr Relationsgraph ist $R|_B=R\cap(B\times B)$.
\end{definition}

\begin{theorem}[Restriktion und wiederholte Restriktion]
\label{b12guide:thm:restriction}
Für jedes $B\subseteq P$ ist $(B,\leq_B)$ eine halbgeordnete Menge.
Für $C\subseteq B\subseteq P$ gilt
\[
 (R|_B)|_C=R|_C,
 \qquad (R|_B)^{\mathrm{op}}=R^{\mathrm{op}}|_B.
\]
Für $T\subseteq B$ sind Minimalität, Maximalität, Minimum und
Maximum in der Teilmenge $T$ unabhängig davon, ob als umgebender
Träger $B$ oder $P$ verwendet wird.
\end{theorem}

\begin{theorem}[Hüllen und untere Mengen in einer Restriktion]
\label{b12guide:thm:relative-closures}
Für $S\subseteq B\subseteq P$ gilt
\[
 \downarrow_B S=B\cap\downarrow_P S,\qquad
 \uparrow_B S=B\cap\uparrow_P S.
\]
Die unteren und oberen Mengen von $B$ sind genau die Spuren
entsprechender Mengen von $P$:
\[
\begin{aligned}
 \Low(B)&=\{B\cap L\mid L\in\Low(P)\},\\
 \Up(B)&=\{B\cap U\mid U\in\Up(P)\}.
\end{aligned}
\]
\end{theorem}

\section{Intervalle und ordnungskonvexe Teilmengen}

\begin{definition}[Abgeschlossene Ordnungsintervalle]
\label{b12guide:def:interval}
Für $a,b\in P$ sei
\[
 [a,b]_P=\{x\in P\mid a\leq x\leq b\}
         =(\uparrow a)\cap(\downarrow b).
\]
Insbesondere wird keine Vergleichbarkeit von $a$ und $b$
vorab vorausgesetzt.
\end{definition}

\begin{theorem}[Intervallgrenzen und Intervallinklusion]
\label{b12guide:thm:interval}
Für $a,b\in P$ gilt
$[a,b]_P\ne\varnothing\Longleftrightarrow a\leq b$.
Ist $a\leq b$, so gilt für beliebige $c,d\in P$
\[
\begin{aligned}
 \min[a,b]_P&=a,\qquad \max[a,b]_P=b,\\
 [a,b]_P\subseteq[c,d]_P
 &\quad\Longleftrightarrow\quad c\leq a\ \land\ b\leq d.
\end{aligned}
\]
Zwei nichtleere abgeschlossene Intervalle sind genau dann gleich,
wenn ihre jeweiligen Endpunkte übereinstimmen.
\end{theorem}

\begin{definition}[Ordnungskonvexität]
\label{b12guide:def:convex}
Eine Teilmenge $C\subseteq P$ heißt \emph{ordnungskonvex}, wenn
\[
 a,b\in C,\ x\in P,\ a\leq x\leq b\quad\Longrightarrow\quad x\in C.
\]
Die \emph{ordnungskonvexe Hülle} $\Conv(S)$ einer Teilmenge $S$
ist der Durchschnitt aller ordnungskonvexen Teilmengen von $P$,
die $S$ enthalten. Diese Familie enthält stets $P$.
\end{definition}

\begin{theorem}[Charakterisierung ordnungskonvexer Mengen]
\label{b12guide:thm:convex-characterization}
Für $C\subseteq P$ sind äquivalent:
\[
\begin{aligned}
 &C\text{ ist ordnungskonvex},\\
 &\forall a,b\in C\;([a,b]_P\subseteq C),\\
 &C=(\uparrow C)\cap(\downarrow C),\\
 &C=U\cap L\text{ für eine obere Menge }U
                    \text{ und eine untere Menge }L.
\end{aligned}
\]
\end{theorem}

\begin{theorem}[Abgeschlossenheit konvexer Mengen]
\label{b12guide:thm:convex-intersections}
Jede untere Menge, jede obere Menge, jedes Intervall $[a,b]_P$
und jede Antikette ist ordnungskonvex. Beliebige Durchschnitte
ordnungskonvexer Mengen sind ordnungskonvex; der leere
Durchschnitt ist $P$. Für $B\subseteq P$ und ordnungskonvexes
$C\subseteq P$ ist $C\cap B$ ordnungskonvex in der induzierten
Ordnung auf $B$.
\end{theorem}

\begin{theorem}[Explizite konvexe Hülle]
\label{b12guide:thm:convex-hull}
Für jede Teilmenge $S\subseteq P$ gilt
\[
 \Conv(S)=(\uparrow S)\cap(\downarrow S)
          =\bigcup_{a,b\in S}[a,b]_P.
\]
Dies ist die kleinste ordnungskonvexe Menge, die $S$ enthält.
Insbesondere gelten
\[
\begin{aligned}
 S&\subseteq\Conv(S),& \Conv(\Conv(S))&=\Conv(S),\\
 S\subseteq T&\ \Rightarrow\ \Conv(S)\subseteq\Conv(T),&
 \Conv(\varnothing)&=\varnothing.
\end{aligned}
\]
\end{theorem}
